\documentclass{article}
\usepackage[T1]{fontenc}
\usepackage{lmodern}
\usepackage{graphicx}
\usepackage{amsmath,amssymb}
\usepackage[a4paper,margin=2cm]{geometry}
\usepackage[absolute,overlay]{textpos}
\setlength{\TPHorizModule}{1mm}
\setlength{\TPVertModule}{1mm}
\usepackage[indonesian]{babel}
\usepackage{hyperref}
\setlength{\emergencystretch}{2em}
% unit: Halaman 2

\begin{document}

% === Halaman 2 (python/native) ===

2

Pendahuluan 

• Keseluruhan point dari kriptografi adalah menjaga kerahasiaan pesan 

atau kunci dari pihak ketiga. Pihak ketiga: musuh, lawan, penyadap, 

kriptanalis (cryptanalyst). 

• Pihak lawan berusaha memecahkan cipherteks dengan melakukan 

serangan terhadap sistem kriptografi.

• Tujuan serangan adalah untuk mengungkap plainteks dari cipherteks 

atau mendapatkan kunci dekripsi, sehingga kunci dekripsi dapat 

digunakan untuk mendekripsi potongan cipherteks lainnya.

\end{document}
